% !PRINTED 47
\weberpage{47}
%<<CONT
und indem wir in dieser identischen Gleichung $x = 0, 1, 1:\varkappa^2$
setzen, so folgt:
\weq{16}{\begin{aligned}
\sqrt{x_1x_2x_3} &= -\dfrac{c}{b} \\
\sqrt{(1-x_1)(1-x_2)(1-x_3)} &= \dfrac{b+a}{b} \\
\sqrt{(1-\varkappa^2x_1)(1-\varkappa^2x_2)(1-\varkappa^2x_3)} &= \dfrac{b+\varkappa^2a}{b}.
\end{aligned}}

Zur Bestimmung der Vorzeichen in diesen Gleichungen erhalten
 wir noch aus (13), wenn wir $x_1$ durch $x_2$ und $x_3$ ersetzen
und die Ergebnisse multiplizieren:
\[ x_1x_2x_3(a+bx_1)(a+bx_2)(a+bx_3) = -c^3\sqrt{f(x_1)}\sqrt{f(x_2)}\sqrt{f(x_3)} \]
und (15) ergibt für $x = -a:b$
\[ (a+bx_1)(a+bx_2)(a+bx_3) = \dfrac{c^2}{b}(b+a)(b+\varkappa^2a). \]
Benutzt man noch das Quadrat der ersten Gleichung (16),
so folgt hieraus
\[ \sqrt{f(x_1)}\sqrt{f(x_2)}\sqrt{f(x_3)} = -\dfrac{c}{b^3}(b+a)(b+\varkappa^2a). \]

Dasselbe Resultat ergibt aber die Multiplikation der drei
Gleichungen (16), und die Vorzeichen sind also in diesen Formeln
so gewählt, daß das Produkt der linken Seiten den durch die
Differentialgleichung (1) geforderten Wert von $\sqrt{f(x_1)}$ ergibt.
Eine weitere Bestimmung der Vorzeichen ist in (16) der Natur
der Sache nach nicht möglich, und diese Gleichungen dienen zur
eindeutigen Definition von $\sqrt{x_1}, \sqrt{1-x_1}, \sqrt{1-\varkappa^2x_1}$, wenn die
$\sqrt{x_2}, \sqrt{1-x_2}, \sqrt{1-\varkappa^2x_2}, \sqrt{x_3}, \sqrt{1-x_3}, \sqrt{1-\varkappa^2x_3}$ als gegeben
vorausgesetzt werden.

Um die Gleichungen (16) in die gebräuchliche Form zu
bringen, machen wir zunächst in (14) den Zähler von $b$ rational
und erhalten
\[ b = -\dfrac{(x_2-x_3)(1-\varkappa^2x_2x_3)}{\sqrt{x_2}\sqrt{(1-x_3)(1-\varkappa^2x_3)} + \sqrt{x_3}\sqrt{(1-x_2)(1-\varkappa^2x_2)}}, \]
ferner nach (14):
\[ \begin{aligned}
b+a &= -\sqrt{(1-x_2)(1-x_3)}\left\{\sqrt{x_2(1-x_2)}\sqrt{1-\varkappa^2x_3} - \sqrt{x_3(1-x_3)}\sqrt{1-\varkappa^2x_2}\right\} \\
b+\varkappa^2a &= -\sqrt{(1-\varkappa^2x_2)(1-\varkappa^2x_3)}\left\{\sqrt{x_2(1-\varkappa^2x_2)}\sqrt{1-x_3} - \sqrt{x_3(1-\varkappa^2x_3)}\sqrt{1-x_2}\right\},
\end{aligned} \]
%CONT>>
